% ==============================================================================
% CHAPITRE 4 : IMPLÉMENTATION DU PIPELINE STREAMING & DÉTECTION
% ==============================================================================

\chapter{Implémentation du Pipeline Streaming et Détection}
\label{ch:streaming}

\section{Validation sémantique et Machine à États Finis (FSM)}

Avant injection dans le graphe actif, le module de validation applique cinq contrôles déterministes modélisés par une Machine à États Finis (FSM, Figure~\ref{fig:fsm_validation}) : conformité du schéma (8 champs JSON obligatoires, montants positifs), devise \texttt{MAD}, exclusion des auto-virements $u_{src} = u_{tgt}$ (test AT-05), contrôle d'horodatage futur $\tau_{event} \le t_{now} + 60\text{ s}$ (test AT-14) et détection des conflits par hachage de la charge utile $H_{payload} = \operatorname{SHA-256}(u_{src} \mid u_{tgt} \mid m \mid \text{devise} \mid \tau)$. Si le même \texttt{transaction\_id} réapparaît avec des attributs altérés, il est consigné dans \texttt{conflict\_transactions} (test AT-12).

\singlefigure[2.1cm]{figures/fig_fsm_validation_states.png}{Machine à états finis (FSM) régissant la validation sémantique et l'admission des transactions}{fig:fsm_validation}

\section{Extraction de motifs et invariance par rotation (AT-02)}

À chaque micro-lot (déclenché toutes les 5 secondes), Spark filtre les transactions de la tranche $[H - 600\text{ s}, H]$, instancie un objet \texttt{GraphFrame} et exécute la recherche de motifs :
\begin{equation}
    \texttt{g.find("(a)-[e1]->(b); (b)-[e2]->(c); (c)-[e3]->(a)")}
\end{equation}
suivie d'un filtre relationnel imposant $a \ne b \ne c$, la causalité $e_1.ts \le e_2.ts \le e_3.ts$ et le ratio de conservation $\frac{\max(m_i)}{\min(m_i)} \le 1.15$. Pour éviter de générer trois alertes distinctes pour les trois permutations rotatoires $(e_1, e_2, e_3)$, $(e_2, e_3, e_1)$ et $(e_3, e_1, e_2)$, le système trie lexicographiquement les identifiants de transactions avant de calculer les empreintes :
\begin{equation}
    \cycleID = \operatorname{SHA-256}\big(\operatorname{Concat}(\operatorname{sort}([txn_1, txn_2, txn_3]))\big), \quad \alertID = \operatorname{SHA-256}\big(\cycleID \mid \text{rule\_id} \mid \text{version}\big)
\end{equation}
Cette normalisation garantit une **invariance stricte par rotation** (test AT-02 validé).

\screenshotpaired[2.5cm]%
{figures/fig_screenshot_spark_ui.png}{Console Spark Web UI (Plan d'exécution DAG)}%
{figures/fig_screenshot_kafka_console.png}{Console Kafka UI (Topic et partitionnement)}%
{Supervision opérationnelle en temps réel de l'infrastructure de streaming sous charge}%
{fig:infra_telemetry}

\section{Sauvegarde Delta, publication Kafka et isolation des checkpoints}

L'alerte validée est insérée dans la table Delta Lake \texttt{fraud\_alerts} via une clause \texttt{MERGE INTO ... WHEN NOT MATCHED THEN INSERT *}, assurant l'idempotence totale d'écriture. L'alerte est simultanément publiée sur le topic Kafka \texttt{fraud-alerts} avec l'\alertID{} comme clé de déduplication pour les consommateurs. Enfin, pour prévenir les corruptions d'état en cas de redémarrage après arrêt ou crash (test AT-19), le système isole rigoureusement le répertoire de streaming (\texttt{data/checkpoints/}) du répertoire de linéarité GraphFrames (\texttt{data/graph\_checkpoints/}).

\newpage
